\section{Extended Mathematical Details}
\label{app:extended-mathematics}
The core queue, pressure, imbalance, and reward definitions appear in the main
paper. This section retains their reference-scale calibration, the offline
objective, and the full proofs.
\subsection{Variance Reference Scales}
The subscription-count reference $V^N_{\mathrm{ref}}$ is chosen from an
operational tolerance rather than learned by DQL. Since
$\sqrt{V^N_u(t)}$ is the standard deviation of subscription counts across
client $u$'s active connections, a tolerable subscription-count standard
deviation $\sigma^{N}_{\mathrm{tol}}$ defines
\[
	V^N_{\mathrm{ref}}=(\sigma^{N}_{\mathrm{tol}})^2.
\]
For example, if the operator allows roughly a $\Delta_N$-subscription spread
between a heavy and a light connection, then
$\sigma^{N}_{\mathrm{tol}}\approx \Delta_N/2$ is a practical calibration. If
the subscription-density threshold is $\Theta_N$, a default choice is
$\sigma^{N}_{\mathrm{tol}}=\rho_N\Theta_N$ with
$\rho_N\in[0.1,0.25]$; smaller values increase sensitivity to
subscription-count skew.
The backlog-load reference $V^G_{\mathrm{ref}}$ is chosen in the same way,
but in backlog units. Since $\sqrt{V^G_u(t)}$ is the standard deviation of
cumulative connection backlog $G_{u,c}(t)$ across client $u$'s active
connections, a tolerable backlog-load standard deviation
$\sigma^{G}_{\mathrm{tol}}$ defines
\[
	V^G_{\mathrm{ref}}=(\sigma^{G}_{\mathrm{tol}})^2.
\]
If the operator allows roughly a $\Delta_G$-message backlog spread between a
heavy and a light connection, then
$\sigma^{G}_{\mathrm{tol}}\approx \Delta_G/2$ is the direct calibration. When
the tolerance is latency-driven, $\Delta_G$ is derived from the
client-observed processing capability rather than from pod capacity alone. For
example, if $\widehat{r}^{\mathrm{ack}}_{u}$ is a conservative ACK or processing rate
measured by client $u$ and $\tau_G$ is the acceptable additional drain-time
skew, then $\Delta_G=\widehat{r}^{\mathrm{ack}}_{u}\tau_G$ is a defensible choice. If
no latency SLO is available, a practical default is
$\sigma^{G}_{\mathrm{tol}}=\rho_G\beta\Theta_N$ with
$\rho_G\in[0.1,0.25]$, where $\beta\Theta_N$ is the nominal backlog budget of
a connection carrying $\Theta_N$ subscriptions.
The weights $\eta_N,\eta_G\geq0$ determine how much the controller prioritizes
subscription-management fairness versus backlog-load fairness. For example,
larger $\eta_N$ makes the controller more sensitive to connections carrying
too many subscriptions, while larger $\eta_G$ makes it more sensitive to
connections accumulating more queued messages. The constraint
$\eta_N+\eta_G=1$ makes $V_u(t)$ a convex combination rather than an
unbounded sum.
\subsection{Global Evaluation Objective}

The controller does not require cluster-wide variance. However, simulation and
offline validation can report global metrics. A finite-horizon evaluation cost
is
\begin{align}
	J_T =
	\sum_{t=0}^{T}
	\Bigl[
	&\Lambda_B \sum_{i\in\mathcal{S}}B_i(t)
	+\Lambda_X \Phi_X(t) \nonumber\\
	&+\Lambda_C \sum_{u\in\mathcal{U}}V_u(t)
	+\Lambda_M M(t)
	\Bigr].
\end{align}
where $\Phi_X(t)$ is pod subscription-count variance and $M(t)$ is total
migration count. This objective is used to evaluate system-level behavior, not
as information assumed available to every client.

The quantity $J_T$ is an offline finite-horizon cost: lower values indicate
better system-level behavior over the evaluation horizon $0,\ldots,T$. The
first term penalizes total pod backlog, where $B_i(t)$ is the cumulative
backlog over all subscription queues served by pod $i$. The second term
penalizes uneven subscription placement across pods. Specifically, if
\begin{align*}
	\bar{X}(t)
	&=\frac{1}{|\mathcal{S}|}\sum_{i\in\mathcal{S}}X_i(t),\\
	\Phi_X(t)
	&=\frac{1}{|\mathcal{S}|}
	\sum_{i\in\mathcal{S}}\bigl(X_i(t)-\bar{X}(t)\bigr)^2,
\end{align*}
then $\Phi_X(t)$ is high when some pods carry many more subscriptions than
others. The third term aggregates client-side connection imbalance, using
$V_u(t)$ from (\mainref{eq:local_connection_imbalance}); it measures whether a
client's own persistent connections are uneven in subscription count or
backlog load. The fourth term penalizes operational churn through $M(t)$, the
number of migration or reconnect operations performed during epoch $t$.
The coefficients
$\Lambda_B,\Lambda_X,\Lambda_C,\Lambda_M\geq 0$ determine the relative
importance of backlog reduction, pod-level subscription balance, client-side
connection balance, and migration stability. Normalization before weighting
places the components on comparable scales, preventing a large-scale metric
from dominating $J_T$ merely because of its units.

\subsection{Rebalancing Properties}

\begin{proposition}[Subscription-density bound]
If $1\leq|\mathcal{K}_u(t)|\leq \Theta_N C_{\max}$ and the scaling rule sets
\begin{equation}
	|\mathcal{C}_u(t)| =
	\left\lceil\frac{|\mathcal{K}_u(t)|}{\Theta_N}\right\rceil,
\end{equation}
then the mean subscription density satisfies $\bar{N}_u(t)\leq\Theta_N$.
\end{proposition}

\begin{proof}
By definition, $\bar{N}_u(t)=|\mathcal{K}_u(t)|/|\mathcal{C}_u(t)|$. Since
$|\mathcal{C}_u(t)|\geq |\mathcal{K}_u(t)|/\Theta_N$, it follows that
$\bar{N}_u(t)\leq\Theta_N$.
\end{proof}

\begin{proposition}[Backlog-variance reduction]
For client $u$, let $c_h,c_l\in\mathcal{C}_u(t)$ be two active persistent
connections, where $c_h$ is the source connection from which subscription $k$
is removed and $c_l$ is the target connection to which subscription $k$ is
assigned. Let $G_{u,c_h}(t)$ and $G_{u,c_l}(t)$ be the cumulative backlog loads
on these two connections before the move, and let $b_{u,k}(t)>0$ be the backlog
of the moved subscription queue. If
\begin{equation}
	G_{u,c_h}(t)-G_{u,c_l}(t)>b_{u,k}(t),
\end{equation}
then moving subscription $k$ from $c_h$ to $c_l$ strictly decreases the
connection-level backlog variance $V^G_u(t)$.
\end{proposition}

The condition states that the source connection is heavier than the target
connection by more than the backlog carried by the moved subscription. Hence
the move reduces their squared-deviation sum. It can cross the equal-load
point and still reduce variance: loads $(10,0)$ become $(2,8)$ after a move
of size 8. Avoiding that crossing would require the stronger condition
$2b_{u,k}(t)\leq G_{u,c_h}(t)-G_{u,c_l}(t)$, which is not needed by the proposition.

\begin{proof}
The move preserves total backlog and therefore preserves $\bar{G}_u(t)$. Let
$d_h=G_{u,c_h}(t)-\bar{G}_u(t)$ and
$d_l=G_{u,c_l}(t)-\bar{G}_u(t)$. Only two squared-deviation terms change. With
$b=b_{u,k}(t)$, the unnormalized variance change is
\begin{align}
	&(d_h-b)^2+(d_l+b)^2-d_h^2-d_l^2 \nonumber\\
	&\quad =
	2b\bigl(b-(G_{u,c_h}(t)-G_{u,c_l}(t))\bigr)<0.
\end{align}
Thus $V^G_u(t)$ decreases.
\end{proof}

\begin{remark}
The placement loop optimizes a client-observed reward rather than an
unobserved global cluster objective. It is a decentralized control mechanism operating
under partial observability: active connections and probe connections provide
fresh but incomplete pod statistics, and the client optimizes its own reward
using these observations.
\end{remark}

